% Source: https://artofproblemsolving.com/community/c3h2310037_challenging_geometry_problem
% Stable identity: aops-topic-2310037-q01
Label the points so that $A$, $B$, $C$, $D$ lie on the circle in this order.

\begin{asy}
usepackage("amsmath");
size(5.5cm);
defaultpen(fontsize(10pt));
pair O=(0,0);
pair A=(-sqrt(0.96),-0.2), C=(sqrt(0.96),-0.2), B=(0.3,sqrt(0.91)), D=(0.3,-sqrt(0.91));
pair P=(0.3,-0.2);
draw(circle(O,1),gray);
draw(A--C,linewidth(1)); draw(B--D,linewidth(1));
draw(A--B,blue); draw(C--D,red); draw(B--C,dashed);
draw(P+(0.08,0)--P+(0.08,0.08)--P+(0,0.08));
draw(O--dir(140),gray+dashed); label("$R$",0.5*dir(140)+0.06*dir(230),dir(230),gray);
dot(O); dot(P);
label("$O$",O,N); label("$P$",P,SE);
label("$A$",A,W); label("$B$",B,N); label("$C$",C,E); label("$D$",D,S);
\end{asy}

\textbf{(a)} In triangles $APB$ and $DPC$, the angles at $P$ are vertical angles, and $\angle BAC=\angle BDC$ (inscribed angles on the arc $BC$). So the triangles are similar, and $\dfrac{AP}{DP}=\dfrac{BP}{CP}$, i.e. $AP\cdot CP=BP\cdot DP$.

\textbf{(b)} Now $\angle APB=\angle CPD=90^\circ$, so by Pythagoras
\[AP^2+BP^2=AB^2,\qquad CP^2+DP^2=CD^2.\]
The chord $AB$ subtends the inscribed angle $\angle ACB$ and the chord $CD$ subtends the inscribed angle $\angle CBD$, so by the extended law of sines
\[AB=2R\sin\angle ACB,\qquad CD=2R\sin\angle CBD.\]
In triangle $PBC$ the angle at $P$ is $90^\circ$, so $\angle PCB+\angle PBC=90^\circ$, that is, $\angle ACB+\angle CBD=90^\circ$, and $\sin\angle CBD=\cos\angle ACB$. Therefore
\[AP^2+BP^2+CP^2+DP^2=AB^2+CD^2=4R^2\left(\sin^2\angle ACB+\cos^2\angle ACB\right)=4R^2.\]
$\blacksquare$

(Without the perpendicularity there is no relation between the four lengths and $R$ other than (a): keeping $AP,BP,CP,DP$ fixed and changing the angle between the chords changes $R$, unless all four lengths are equal and $P$ is the centre.)
